\documentclass[12pt]{article}
\usepackage{amsmath,amssymb,amsfonts}
\usepackage[margin=1in]{geometry}

\begin{document}

\textbf{Chapter 6, Problem 20.} Prove that if $f(z)$ is analytic within and on a closed contour $C$, and $a$ and $b$ are two distinct points within $C$, then
\[
\frac{1}{2\pi i}\oint_C \frac{f(z)}{(z-a)(z-b)}\,dz = \frac{f(a)-f(b)}{a-b}.
\]

From this result, deduce Liouville's theorem, which states that a function must be a constant if it is entire and bounded: $|f(z)| \leq M$ for all $z$.

\bigskip
\textbf{Solution.}

\textbf{Part 1: Proof of the identity.}

Use partial fractions:
\[
\frac{1}{(z-a)(z-b)} = \frac{1}{a-b}\left(\frac{1}{z-a} - \frac{1}{z-b}\right).
\]

Therefore:
\[
\frac{1}{2\pi i}\oint_C \frac{f(z)}{(z-a)(z-b)}\,dz = \frac{1}{a-b}\left[\frac{1}{2\pi i}\oint_C\frac{f(z)}{z-a}\,dz - \frac{1}{2\pi i}\oint_C\frac{f(z)}{z-b}\,dz\right].
\]

By Cauchy's integral formula (since $f$ is analytic inside $C$ and both $a,b$ are inside $C$):
\[
\frac{1}{2\pi i}\oint_C\frac{f(z)}{z-a}\,dz = f(a), \qquad \frac{1}{2\pi i}\oint_C\frac{f(z)}{z-b}\,dz = f(b).
\]

Therefore:
\[
\frac{1}{2\pi i}\oint_C \frac{f(z)}{(z-a)(z-b)}\,dz = \frac{f(a)-f(b)}{a-b}. \quad \checkmark
\]

\textbf{Part 2: Liouville's theorem.}

Suppose $f$ is entire and $|f(z)| \leq M$ for all $z$. Let $C$ be a circle of radius $R$ centered at the origin, with $R$ large enough to enclose both $a$ and $b$.

On $C$: $|z| = R$, $|z-a| \geq R - |a|$, $|z-b| \geq R - |b|$. By the ML inequality:
\[
\left|\frac{f(a)-f(b)}{a-b}\right| = \left|\frac{1}{2\pi i}\oint_C \frac{f(z)}{(z-a)(z-b)}\,dz\right| \leq \frac{1}{2\pi}\cdot\frac{M}{(R-|a|)(R-|b|)}\cdot 2\pi R = \frac{MR}{(R-|a|)(R-|b|)}.
\]

As $R \to \infty$:
\[
\frac{MR}{(R-|a|)(R-|b|)} \to \frac{M}{R} \to 0.
\]

Therefore $\frac{f(a)-f(b)}{a-b} = 0$, which gives $f(a) = f(b)$ for all $a, b$. Hence $f$ is constant.

\end{document}
